\documentclass{amsart}
% For dealing with references we use the comment environment
\usepackage{verbatim}
\newenvironment{reference}{\comment}{\endcomment}
%\newenvironment{reference}{}{}
\newenvironment{slogan}{\comment}{\endcomment}
\newenvironment{history}{\comment}{\endcomment}

% For commutative diagrams we use Xy-pic
\usepackage[all]{xy}

% We use 2cell for 2-commutative diagrams.
\xyoption{2cell}
\UseAllTwocells

% We use multicol for the list of chapters between chapters
\usepackage{multicol}

% This is generally recommended for better output
\usepackage{lmodern}
\usepackage[T1]{fontenc}

% For cross-file-references

% Package for hypertext links:
\usepackage{hyperref}

% For any local file, say "hello.tex" you want to link to please

% Theorem environments.
%
\theoremstyle{plain}
\newtheorem{theorem}[subsection]{Theorem}
\newtheorem{proposition}[subsection]{Proposition}
\newtheorem{lemma}[subsection]{Lemma}

\theoremstyle{definition}
\newtheorem{definition}[subsection]{Definition}
\newtheorem{example}[subsection]{Example}
\newtheorem{exercise}[subsection]{Exercise}
\newtheorem{situation}[subsection]{Situation}

\theoremstyle{remark}
\newtheorem{remark}[subsection]{Remark}
\newtheorem{remarks}[subsection]{Remarks}

\numberwithin{equation}{subsection}

% Macros
%
\def\lim{\mathop{\mathrm{lim}}\nolimits}
\def\colim{\mathop{\mathrm{colim}}\nolimits}
\def\Spec{\mathop{\mathrm{Spec}}}
\def\Hom{\mathop{\mathrm{Hom}}\nolimits}
\def\Ext{\mathop{\mathrm{Ext}}\nolimits}
\def\SheafHom{\mathop{\mathcal{H}\!\mathit{om}}\nolimits}
\def\SheafExt{\mathop{\mathcal{E}\!\mathit{xt}}\nolimits}
\def\Sch{\mathit{Sch}}
\def\Mor{\mathop{\mathrm{Mor}}\nolimits}
\def\Ob{\mathop{\mathrm{Ob}}\nolimits}
\def\Sh{\mathop{\mathit{Sh}}\nolimits}
\def\NL{\mathop{N\!L}\nolimits}
\def\CH{\mathop{\mathrm{CH}}\nolimits}
\def\proetale{{pro\text{-}\acute{e}tale}}
\def\etale{{\acute{e}tale}}
\def\QCoh{\mathit{QCoh}}
\def\Ker{\mathop{\mathrm{Ker}}}
\def\Im{\mathop{\mathrm{Im}}}
\def\Coker{\mathop{\mathrm{Coker}}}
\def\Coim{\mathop{\mathrm{Coim}}}

% Boxtimes
%
\DeclareMathSymbol{\boxtimes}{\mathbin}{AMSa}{"02}

%
% Macros for moduli stacks/spaces
%
\def\QCohstack{\mathcal{QC}\!\mathit{oh}}
\def\Cohstack{\mathcal{C}\!\mathit{oh}}
\def\Spacesstack{\mathcal{S}\!\mathit{paces}}
\def\Quotfunctor{\mathrm{Quot}}
\def\Hilbfunctor{\mathrm{Hilb}}
\def\Curvesstack{\mathcal{C}\!\mathit{urves}}
\def\Polarizedstack{\mathcal{P}\!\mathit{olarized}}
\def\Complexesstack{\mathcal{C}\!\mathit{omplexes}}
% \Pic is the operator that assigns to X its picard group, usage \Pic(X)
% \Picardstack_{X/B} denotes the Picard stack of X over B
% \Picardfunctor_{X/B} denotes the Picard functor of X over B
\def\Pic{\mathop{\mathrm{Pic}}\nolimits}
\def\Picardstack{\mathcal{P}\!\mathit{ic}}
\def\Picardfunctor{\mathrm{Pic}}
\def\Deformationcategory{\mathcal{D}\!\mathit{ef}}

\setlength{\emergencystretch}{3em}
\begin{document}
\title[Stacks Project, Tag 04PZ]{576 --- Connected components map onto connected components under arbitrary field extension for arbitrary schemes}
\maketitle
\noindent Copyright (C) 2005--2025 Johan de Jong. Stacks Project authors. Source: \href{https://stacks.math.columbia.edu/tag/04PZ}{Tag 04PZ}.

\noindent Editorial difficulty note (not author text). Difficulty H3: The proof cannot rely on finite type or quasi-compactness of the scheme. It isolates the maximal weakly étale subfield of global functions, proves geometric connectedness over that field, and identifies every extended component as a whole fibre over a field quotient of the tensor product..

\noindent Editorial provenance note (not author text). History: excerpt from revision \texttt{a0444\allowbreak 6e57e\allowbreak c1fbc\allowbreak 252a8\allowbreak 71afc\allowbreak ec775\allowbreak 2fb28\allowbreak 07b14}, varieties.tex, lines 1147--1223. Reference presentation was adapted; original-author context and core support blocks were retained or added. The mathematical wording is unchanged. Licensed under GNU FDL 1.2 or later; no invariant sections or cover texts. See ../licenses/COPYING. Context blocks reproduce author definitions and supporting results; they are not additional counted problems.

\begin{lemma}[Gabber]
\label{lemma-image-connected-component}
\begin{reference}
Email from Ofer Gabber dated June 4, 2016
\end{reference}
Let $K/k$ be an extension of fields. Let $X$ be a scheme over $k$.
Denote $p : X_K \to X$ the projection morphism.
Let $\overline{T} \subset X_K$ be a connected component.
Then $p(\overline{T})$ is a connected component of $X$.
\end{lemma}

\begin{proof}
When $K/k$ is finite this is
Lemma \href{https://stacks.math.columbia.edu/tag/07VM}{lemma image connected component finite extension (Tag 07VM)}.
In general the proof is more difficult.

\medskip\noindent
Let $T \subset X$ be the connected component of $X$ containing
the image of $\overline{T}$. We may replace $X$ by $T$
(with the induced reduced subscheme structure). Thus we
may assume $X$ is connected. Let $A = H^0(X, \mathcal{O}_X)$.
Let $L \subset A$ be the maximal weakly \'etale $k$-subalgebra, see
More on Algebra, Lemma \href{https://stacks.math.columbia.edu/tag/0CKS}{lemma max weakly etale subalgebra (Tag 0CKS)}.
Since $A$ does not have any nontrivial idempotents we see
that $L$ is a field and a separable algebraic extension of $k$ by
More on Algebra, Lemma \href{https://stacks.math.columbia.edu/tag/0CKR}{lemma class weakly etale over field (Tag 0CKR)}.
Observe that $L$ is also the maximal weakly \'etale $L$-subalgebra of $A$
(because any weakly \'etale $L$-algebra is weakly \'etale over $k$
by More on Algebra, Lemma \href{https://stacks.math.columbia.edu/tag/092J}{lemma composition weakly etale (Tag 092J)}).
By Schemes, Lemma \href{https://stacks.math.columbia.edu/tag/01I1}{lemma morphism into affine (Tag 01I1)}
we obtain a factorization $X \to \Spec(L) \to \Spec(k)$
of the structure morphism.

\medskip\noindent
Let $L'/L$ be a finite separable extension. By
Cohomology of Schemes, Lemma
\href{https://stacks.math.columbia.edu/tag/0CKW}{lemma finite locally free base change cohomology (Tag 0CKW)}
we have
$$
A \otimes_L L' =
H^0(X \times_{\Spec(L)} \Spec(L'), \mathcal{O}_{X \times_{\Spec(L)} \Spec(L')})
$$
The maximal weakly \'etale $L'$-subalgebra of $A \otimes_L L'$
is $L \otimes_L L' = L'$ by More on Algebra, Lemma
\href{https://stacks.math.columbia.edu/tag/0CKU}{lemma change fields max weakly etale subalgebra (Tag 0CKU)}.
In particular $A \otimes_L L'$ does not have nontrivial idempotents
(such an idempotent would generate a weakly \'etale subalgebra)
and we conclude that $X \times_{\Spec(L)} \Spec(L')$ is connected.
By Lemma \href{https://stacks.math.columbia.edu/tag/0389}{lemma characterize geometrically disconnected (Tag 0389)}
we conclude that $X$ is geometrically connected over $L$.

\medskip\noindent
Let's give $\overline{T}$ the reduced induced scheme structure
and consider the composition
$$
\overline{T} \xrightarrow{i} X_K = X \times_{\Spec(k)} \Spec(K)
\xrightarrow{\pi}
\Spec(L \otimes_k K)
$$
The image is contained in a connected component of $\Spec(L \otimes_k K)$.
Since $K \to L \otimes_k K$ is integral we see that
the connected components of $\Spec(L \otimes_k K)$
are points and all points are closed, see
Algebra, Lemma \href{https://stacks.math.columbia.edu/tag/00GS}{lemma integral over field (Tag 00GS)}.
Thus we get a quotient field $L \otimes_k K \to E$
such that $\overline{T}$ maps into $\Spec(E) \subset \Spec(L \otimes_k K)$.
Hence $i(\overline{T}) \subset \pi^{-1}(\Spec(E))$. But
$$
\pi^{-1}(\Spec(E)) =
(X \times_{\Spec(k)} \Spec(K)) \times_{\Spec(L \otimes_k K)} \Spec(E) =
X \times_{\Spec(L)} \Spec(E)
$$
which is connected because $X$ is geometrically connected over $L$.
Then we get the equality
$\overline{T} = X \times_{\Spec(L)} \Spec(E)$ (set theoretically)
and we conclude that $\overline{T} \to X$ is surjective as desired.
\end{proof}
\end{document}
